\subsection{The pure source of an arbitrary strict expression}
\label{ca:strict-pure-section}

Return to the full strict syntax of Section~\ref{ca:intro}, with
$b\ge2$, and set every retained sign to zero in the source.
This does not yet change a filling: Section~\ref{ca:filling} will
expand each original full power into these pure circle operators.
Let $y_*$ be the strict complete maximizer and let $a_\ell$ be its
leaf coordinates. Choose a rectangular window $\mathcal B$ about
$y_*$ in the primitive coordinates of \cite[Lemma 3.1]{Profiles}.
All its leaf coordinates are in a compact subset of $(0,1/2)$.
There are still exactly $k=2b-3$ free coordinates and $b-2$ retained
products, with the same integer offsets \eqref{ca:offsets}.
This includes $b=2$, when there is one free coordinate and no
retained product.

For integer wires the pure leaf is exactly \eqref{ca:pure-leaf};
that formula is valid below as well as above $r=N/4$.
Choose the color disk so small that every $r/N\pm t$ in the
window stays in the factorial strip. Define $C_N(t)$ by the
finite window sum of these pure leaves and \eqref{ca:gate}.
For real $t$ near zero it is positive for every sufficiently
large $N$, and it is separately even and holomorphic.
For complex wires its exact interpolating leaf is
\begin{equation}\label{ca:pure-index}
 \mathscr L_N(x,t)=
 \frac{N^2G_N(2Nx)}
 {G_N(Nx)^2G_N(N(x-t))G_N(N(x+t))}.
\end{equation}
At $x=r/N$ this is the original leaf, with no approximation.
Interpolate every gate by replacing each $S_j$ in
\eqref{ca:gate} by $G_N(j)$. Thus the entire product agrees
exactly on the primitive grid, with every affine wire offset.
There is no recurrence to continue across a quarter index.

\begin{theorem}[The strict pure source]\label{ca:strict-pure}
For every fixed strict expression with $b\ge2$, the finite
function just defined has the uniform expansion
\[
 C_N(t)=N^{3/2}a_0(t)e^{N\Phi_0(t)}(1+O(N^{-1}))
\]
on a fixed complex polydisk, with all fixed derivatives on smaller
polydisks. The functions $a_0,\Phi_0$ are holomorphic and separately
even, $a_0$ is nonzero, $a_0(0)>0$, and
\begin{equation}\label{ca:strict-pure-hessian}
 \Phi_0(0)=F_0,\qquad
 \Phi_0''(0)=-2\pi\diag(X_1,\ldots,X_b),\qquad
 X_\ell=\cot(\pi a_\ell)>0.
\end{equation}
The original physical pure source satisfies
\eqref{ca:source-error} with $\epsilon=0$, and has a strict
exponential loss when any color deficit is at least $\delta N$.
The complex Gaussian bound \eqref{ca:color-gaussian} also holds
for this exact $C_N$.
\end{theorem}

\begin{proof}
The phase of the interpolated summand is
\[
 F(y,t)=
 \sum_{\ell=1}^b
 \{J(2L_\ell)-2J(L_\ell)-J(L_\ell-t_\ell)-J(L_\ell+t_\ell)\}
 +\sum_{\text{gates}}m(L_i,L_j,L_r).
\]
All arguments lie in fixed strict complex neighborhoods.
Lemma~\ref{ca:factorial} therefore gives, uniformly with fixed
derivatives, the summand
$N^{3-b}b_0(y,t)e^{NF(y,t)}(1+O(N^{-1}))$.
To specify its normalization, the unshifted leaf density is
\[
 \frac{\sqrt{\sin(2\pi x)}}
 {2^{3/2}\sin(\pi x)
       \sqrt{\sin(\pi(x-t))\sin(\pi(x+t))}},
\]
and the gate density is $b_\mu$ displayed in the proof of
Theorem~\ref{ca:packet-theorem}. Multiply these densities, then
multiply by $e^{\delta_\ell g_x}$ at each leaf and by
$e^{\delta_i m_x+\delta_jm_y+\delta_rm_z}$ at each gate.
Here $g(x,t)=J(2x)-2J(x)-J(x-t)-J(x+t)$, and all derivatives
are evaluated at their actual changed arguments.
This defines $b_0$ without an unspecified normalization.
At real strict arguments its complete value is positive.

Theorem 3.3 of \cite{Profiles} gives the negative definite
matrix $F_{yy}(y_*,0)$. The holomorphic implicit-function theorem
gives $y_*(t)$ and
\[
 \Phi_0(t)=F(y_*(t),t),\qquad
 a_0(t)=
 \frac{(2\pi)^{k/2}b_0(y_*(t),t)}
 {\sqrt{\det[-F_{yy}(y_*(t),t)]}}.
\]
Both are independent of $N$. The square root is positive at
zero and is then continued holomorphically. Primitivity of the
coordinates shows that there is no additional lattice covolume.

We give the contour justification because exact interpolation
is the reason this proof has the larger scope.
Move each endpoint of the wire window by at most $1/N$ to a
half-grid point. The exact kernel $\pi N\cot(\pi Nz)$ encloses
precisely the retained grid points. On a vertical half-grid side
its modulus is at most $\pi N$. On horizontal sides of height
$\pm a$ its constant parts are $\mp\pi N\ii$, with error
$O(Ne^{-2\pi Na})$. The real complete action has a fixed loss
$g_\partial>0$ on the boundary of this window. Since it is real
on real wires, imaginary displacement by $a$ has real-action
cost at most $Ca^2$ there. Choose $a>0$ first so that
$Ca^2<\min(\pi a,g_\partial/4)$, and then shrink the color
disk. Analyticity and evenness in the color give an additional
uniform cost $C\|t\|^2$. Every vertical boundary and every
nonconstant cotangent contribution has a fixed exponential loss.
The two constant horizontal contributions have weight $1/2$
each. Repeating this operation in the $k$ coordinates replaces
the sum by $N^k$ times its integral, with exponentially small
absolute error and no factor $2^k$.

Translate the integral to the parallel contour through $y_*(t)$.
The connecting boundary faces retain the same strict loss.
The positive real part of $-F_{yy}$ gives a uniform Gaussian
on a fixed ball; odd cubic and linear terms integrate to zero,
and the remainder is relative $O(N^{-1})$.
The prefactor is $N^{3-b+k/2}=N^{3/2}$.
Cauchy's formula gives the derivative estimates.
This argument uses only the exact function
\eqref{ca:pure-index}, not the signed interpolation of
Section~\ref{ca:leaf}.

Separate evenness gives $F_{yt}(y,0)=0$.
Moreover $g_{tt}(x,0)=-2\pi\cot(\pi x)$.
The stationary-wire Schur term therefore vanishes at zero,
proving \eqref{ca:strict-pure-hessian}. Positivity of the
complete density proves $a_0(0)>0$ and nonvanishing nearby.
Its negative real-color Hessian and holomorphy give
\eqref{ca:color-gaussian}, after shrinking the disk.

It remains to compare this window with the full physical
source, rather than merely with its local summands.
For a physical deficit $h\ge0$ the pure core leaf is zero for
$r<h$; for $h\le r\le M$ it is the nonnegative factorial
ratio \eqref{ca:pure-leaf}. The uniform endpoint estimate
$\log S_r=NJ(r/N)+O(\log N)$ from
\cite[Lemma 5.2]{Newton} controls this ratio on its entire
support. Its rate is
\[
 g(x,t)=g(x,0)-\Delta(x,t),\qquad
 \Delta(x,t)=J(x-t)+J(x+t)-2J(x)\ge0
 \quad(0\le t\le x\le1/2).
\]
Indeed, for $0<t<x<1/2$,
\[
 \partial_t\Delta(x,t)
 =\log\frac{\sin\pi(x+t)}{\sin\pi(x-t)}>0,
\]
because the difference of the two sines is
$2\cos(\pi x)\sin(\pi t)>0$. Continuity gives the boundary
inequality. At displaced finite endpoints the modulus of
continuity of $J$ is $O(\eta|\log\eta|)$, so the fixed wire
offsets still cost only $O(\log N)$.

The complete action has its unique global maximum strictly
inside $\mathcal B$. Its continuous value on the compact
polytope outside the window is therefore below $F_0$ by a fixed
amount. The nonnegative losses $\Delta$ preserve that gap for
every real color in the chosen corridor. The polynomial number
of states and all $O(\log N)$ errors do not remove the gap.
This proves a uniform bound for the entire pure core outside
the window. Finally \cite[Theorem 1.1]{Bounds}, or
\cite[Corollary 9.1]{Bounds} when $b=2$, bounds the whole
full-minus-core difference and supplies the macroscopic-color
loss. Combining these bounds proves the physical source
assertions at both parities.
\end{proof}

This theorem and the following Fourier argument are the local
analytic inputs to the strict volume theorem. Neither uses the
signed recurrence, its separated-root interval, or a quarter-index
asymptotic.
